\section{Capítulo~\ref{sec:eetypes}. Neuronas como instrumentos}

Los modos de trabajo de las células individuales están medidos directamente, con corriente a través de un electrodo y registro de la tensión; la matemática del integrador y la división en clases son teoría clásica, y que el cerebro use la reconfiguración de modos para calcular sigue siendo una hipótesis.

{\footnotesize
\setlength{\tabcolsep}{3pt}\renewcommand{\arraystretch}{1.15}
\begin{longtable}{>{\raggedright}p{0.5cm}>{\raggedright}p{4.0cm}>{\raggedright}p{5.6cm}>{\raggedright}p{2.3cm}>{\raggedright\arraybackslash}p{1.7cm}}
\toprule
N.º & Afirmación & Experimento y qué se midió & Fuente & Estado \\
\midrule\endfirsthead
\toprule
N.º & Afirmación & Experimento y qué se midió & Fuente & Estado \\
\midrule\endhead
1 & Resonador: una corriente lenta que se opone al cambio de tensión convierte la neurona en un filtro pasa banda con máximo entre unos pocos hercios y decenas de hercios (sección~\ref{sec:resonator}) & En rodajas, se inyectó en neuronas \nt{GLUT} una corriente sinusoidal de frecuencia creciente y se midió la impedancia $|Z(f)|$: máximo en $2$--$5$~Hz a $33^\circ$C (unos $7$~Hz a $38^\circ$C); la resonancia la crean corrientes lentas de potasio y catiónica y la amplifica una corriente persistente de sodio. Una revisión muestra el mismo recurso en muchas clases de células & \cite{HuVervaekeStorm2002,HutcheonYarom2000} & medido \\
2 & La corriente lenta se comporta como una inductancia; junto con la capacidad de la membrana forma un circuito resonante & Se midió la respuesta subumbral de un axón a la corriente y se describió con ecuaciones de conductancia linealizadas; el circuito equivalente contiene una inductancia “fenomenológica” cuyo valor depende de la tensión & \cite{Mauro1970} & teoría \\
3 & Constante de tiempo de un grupo con realimentación $\tau_{\text{ef}}=\tau/(1-w)$~\eqref{eq:perfectint}; con $\tau=0.1$~s y $w=0.995$, $20$~s & Se deduce en el capítulo de la ecuación lineal del grupo; como mecanismo para mantener la posición del ojo se propuso en un trabajo teórico & deducción del capítulo; \cite{Seung1996} & teoría \\
4 & El integrador de posición del ojo mantiene la entrada mediante la red, no mediante una propiedad de la célula aislada & Registro intracelular en peces, en el núcleo que mantiene la posición del ojo: tras un giro rápido, la frecuencia de la neurona cambia en escalón y se mantiene; una corriente breve en una sola célula solo causa un desplazamiento pasajero, y los escalones de tensión persisten también por debajo del umbral: los sostienen las entradas sinápticas & \cite{Aksay2001} & medido \\
5 & El integrador sin fugas necesita un reajuste constante mediante aprendizaje; desajustado, olvida o se satura & Se entrenó a peces con una realimentación visual proporcional a $\pm$ la posición del ojo: el mantenimiento se volvía inestable o con fugas, la constante de tiempo de las neuronas caía más de un orden de magnitud, hasta $<1$~s; la realimentación visual normal restablecía poco a poco la estabilidad & \cite{Major2004} & medido \\
6 & Neurona biestable: una entrada breve activa una meseta duradera y la inhibición la apaga; la propiedad la activa la serotonina (sección~\ref{sec:bistable}) & Registro intracelular de neuronas \nt{ACET} que controlan los músculos, en el gato: una excitación sináptica breve lleva la neurona a una actividad duradera y una inhibición breve la reinicia; en el gato espinal, el estado aparece tras administrar un precursor de la serotonina & \cite{Hounsgaard1988} & medido \\
7 & Dos clases: integradores (la frecuencia crece desde cero) y resonadores (en el umbral, de inmediato una frecuencia finita) & Las clases se deducen de la teoría de bifurcaciones. Experimento: en rodajas de corteza, las neuronas \nt{GLUT} empiezan a disparar con una frecuencia tan baja como se quiera (tipo~1), y las neuronas \nt{GABA} rápidas, con un salto a una frecuencia finita (tipo~2) & \cite{Izhikevich2007,Tateno2004} & medido \\
8 & Una misma neurona es integrador o detector de coincidencias: la inhibición acorta la ventana de suma & En rodajas de hipocampo, la ventana en que las entradas excitadoras se suman hasta el umbral es de menos de $2$~ms; la acorta la inhibición directa que llega justo después de la excitación; en las dendritas, donde la inhibición es más débil, la ventana es más ancha & \cite{Pouille2001} & medido \\
9 & Línea de retardo hecha de axones (tabla~\ref{tab:eetypes}) & En la lechuza común se midieron los retardos en los axones que llegan a los detectores de coincidencias: distintas longitudes de axón dan distintos retardos, y los detectores están sintonizados a la diferencia de tiempos de llegada del sonido & \cite{Carr1990} & medido \\
10 & Marcapasos en la vigilia y célula silenciosa en el sueño & Las neuronas \nt{SERO} disparan de forma casi regular durante el día, se frenan en el sueño lento y casi callan en el sueño REM (más detalles en el capítulo~\ref{sec:sleep}) & \cite{JacobsAzmitia1992,McGintyHarper1976} & medido \\
11 & El instrumento lo fijan los canales iónicos y los canales de difusión que los ajustan (proposición~\ref{prop:eetypes}) & Demostrado en redes pequeñas: las sustancias moduladoras cambian las propiedades propias de las neuronas y la fuerza de las sinapsis, de modo que un mismo circuito produce ritmos distintos & \cite{Marder2012} & medido \\
12 & El cerebro usa la reconfiguración de modos para calcular (proposición~\ref{prop:eetypes}) & No hay ningún experimento directo en el que la reconfiguración del modo de una célula sea necesaria para resolver una tarea; solo hay experimentos en los que los moduladores cambian la salida de un circuito & \cite{Marder2012} & hipótesis \\
\bottomrule
\end{longtable}}
